\section{Discussion and Limitations}\label{sec:discussion}
\textbf{Ground in evidence the attacker cannot edit.} Grounding works because the attacker controls the text but not the prices: edits that push the agent against the market's own evidence are exactly the ones it discounts. Writing misinformation that agrees with the market does not help the attacker, and misinformation that also moves prices leaves grounding's gain intact. The conflict-seeking attacker is the strongest we found: against the full layer it gains about \val{CmpAwareAdaptiveConflictVsStdFull}\pp{} under every conflict rule. A likely reason is that grounding removes at most a fraction $\kappa$ of a signal's deviation from the evidence, so the most extreme contrary edit can still change the decision. Testing this attacker against grounding alone, and defending against it, is the main open problem.

\textbf{Do not let the attacker trigger the defense.} A rule that reacts to text--market conflict lets the attacker decide when it fires, and hold-on-conflict raised net ASR in our experiments. Shrinking the trade by the consistency removes this penalty. Risk rules that key on market state, such as volatility and drawdown, do not give the attacker this lever at all.

\textbf{Measure defenses against placebos and replayed attacks.} In our setting \val{ShareNoiseOfNaive}\% of the flips that naive ASR attributes to the attack also occur with harmless paraphrases. Defense comparisons should share attacker random numbers or replay a recorded attack: before we keyed attacker seeds by injection type, identically configured arms differed by up to 8.6 points of ASR, more than every defense effect we measured except grounding's. Placebo correction has a blind spot of its own: a defense that makes the agent noisier, such as the detection filter, raises the placebo floor and looks better on net ASR. Raw and net ASR should therefore be reported together.

\textbf{Limitations.} (i)~The main agent is a transparent surrogate. We replicate the headline result with a local 7B LLM as the reader, not with frontier models or an LLM that makes the decision itself, and rates may differ there. (ii)~News is template-generated, also in the real-price runs. Real headlines would need an LLM-based attacker, because our edit operators rely on template fields, and the price-impact model is a linear stress test, not a market model. (iii)~The microstructure null result needs real order-flow data. (iv)~The detectability filter is a small feature classifier (training accuracy \val{DetAcc}) standing in for an LLM judge, and stops are evaluated at daily closes. (v)~Each seed evaluates only 60 days, so we rely on paired differences over many seeds, not single runs.
